\documentclass[../main.tex]{subfiles}
\begin{document}
%==========================================TIÊU ĐỀ TẬP========================================
\begin{table}[h]
    \begin{adjustwidth}{-1cm}{}
        \begin{tabular}{>{\centering\arraybackslash}p{6cm}|>{\raggedright\arraybackslash}p{12.5cm}}
            \multirow{5}{6cm}
            % Cột bên trái: ÉPISODE và số tập
            {\rainbowcolor
                {\\[-40pt]
                    \flaregothic\fontsize{20pt}{20pt}\selectfont \centering
                    \color{eptype}JOUR\color{black}\\[12pt]
                    \hbrk\fontsize{72pt}{72pt}\selectfont
                    \color{epnum}7
                    \\[18pt]
                    % \flaregothic\fontsize{14pt}{14pt}\selectfont
                    % \color{epattr}BIENVENUE, \uline{\color{Banpire} Mel} !
                    % \flaregothic\fontsize{18pt}{18pt}\selectfont
                    % \color{epattr}ÉPISODE PERDU\\
                    \unrainbowcolor}
                \color{black}}
             &
            % Dòng thứ nhất: tên tập tiếng Nhật
            {\fotskip\fontsize{18pt}{18pt}\selectfont 720\ruby{時間}{分/ふん}\ruby{無}{む}\ruby{人}{じん}\ruby{島}{とう}\ruby{生}{せい}\ruby{活}{かつ}
            } \\[6pt]
            % Dòng thứ hai: tên tập tiếng Anh
             & {\cafeteria\fontsize{18pt}{18pt}\selectfont Anime-like Deserted Island Situation}                                   \\[4pt]
            % Dòng thứ ba: tên tập tiếng Việt
             & {\vnmsans\fontsize{18pt}{18pt}\selectfont Sinh tồn nơi đảo hoang trong 720 \ruby{giờ}{phút}}                        \\[8pt]
            % Dòng thứ tư: Nhân vật xuất hiện
             & {\sen{\fontsize{14pt}{14pt}\selectfont Track used: Hawaii (from Professional Block 2) (MBB - Beach)}}
            \\[18pt]
            % Dòng cuối cùng: Thời gian diễn ra của tập (thường sẽ là ngày phát hành)
             & {\continuum\fontsize{16pt}{16pt}\selectfont Ngày phát hành: 16/08/2021}                                             \\[18pt]
        \end{tabular}
    \end{adjustwidth}
\end{table}
%=========================================TRUYỆN CHÍNH========================================
\color{black}

Đã tròn một tuần kể từ ngày mà bộ ba nhà Hikawa - Kuon, Kaigou và Suzuki - bị mắc kẹt ở thành phố Cầu Vồng.

Vào một ngày giữa hè, nắng gắt và oi bức như muốn nung chảy cả mặt đất, lẽ ra họ nên ngồi trong phòng điều hòa mát lạnh hoặc chí ít cũng tìm lấy một mái hiên trú thân.

Thế nhưng, ba người họ giờ đây lại đang vật vờ trên một bãi cát trắng toát, bỏng rát, xung quanh chỉ toàn tiếng sóng vỗ và mùi biển mặn, với hằng hà cây dừa bao xung quanh họ.

\img{nj-7a}

Cậu con trai ngồi thừ ra trên bãi cát, mặt vẫn còn nhăn nhó vì choáng váng: ``\vie{\dots Tại sao chúng ta\dots}''

``\vie{\dots lại ở đây chứ?}'' cô chị lớn tiếp lời, bực bội nổi cáu, một tau chống hông, tay còn lại xốc daajy cô em đang nằm bất tỉnh bên cạnh mình.

Suzuki khẽ rên rỉ một tiếng khi được chị gái lay gọi. Cô mở mắt ra, hơi nheo lại vì ánh nắng, rồi ngồi dậy một cách khó nhọc. ``\vie{\dots Onee-sama\dots\ đây là đâu vậy ạ?}''

Kaigou không đáp ngay. Cô liếc mắt về phía Kuon, người đang ôm đầu, mồ hôi chảy dọc thái dương, rồi liếc sang phía bên kia, nơi một bóng người khác đang lồm cồm ngồi dậy từ bãi cát, tóc bạc pha chút tím và đỏ rối tung như tổ quạ.

\chrint

\begin{wrapfigure}[34]{r}{6cm}
    \includegraphics[width=8cm]{../../../../Image/Book?minato_model.png}
\end{wrapfigure}

Không ai khác đó là Fuwa Minato (\ruby{不}{ふ}\ruby{破}{わ}\ruby{湊}{みなと}, \ruby{Phư}{Bất}-\ruby{oa}{Phá}\ \ruby{Mi-na-tô}{Tấu}), một người gia nhập vào Nijisanji vào quý 4 năm 2019. Cậu chàng này tự nhận mình là một người dẫn chương trình tài ba, rất thích hát và chơi game với lối nói chuyện rất tự nhiên và cuốn hút cùng đôi mắt tím truyền cảm của mình. Câu khẩu hiệu của cậu ta là ``Nụ cười của bạn bằng ba bữa ăn hàng ngày'' và ``Xin hãy cho tôi đứng đầu.''

Dù là một người dẫn chương trình, cậu ta lại không thể uống được rượu bia - cũng không hẳn là uống được, mà cơ thể cậu không cho phép làm điều đó, cho tới một năm sau khi đã ấm chỗ ở gia đình Cầu Vồng. Cậu được biết tới là có rất nhiều thứ quái dị mà cậu cho là thú vị trong mỗi buổi phát sóng, tới độ tính cách và cách hành xử của cậu khiến nhiều người nhăn mặt.

Khác với vẻ ngoài nhẹ nhàng thường thấy, kỹ năng chơi game của cậu ta rất cao, nhưng đôi khi, như một cơn co giật, những pha tấu hài của cậu ta lại bất ngờ bùng nổ giữa các buổi stream game.

Minato là một người rất đam mê âm nhạc và ca hát. Cậu bắt đầu học guitar dưới ảnh hưởng của anh trai và từng chơi trong một ban nhạc như một sở thích thời sinh viên. Khi mới debut, cậu ta giấu kín kinh nghiệm chơi nhạc của mình, nhưng đến năm thứ ba hoạt động mới chính thức giới thiệu cây đàn guitar Elk của mình.

Mặc dù giọng điệu của cậu ta có vẻ thô lỗ và hay tấu hài trong các buổi phát sóng, nhưng thực tế Minato lại là người không bao giờ nổi giận với những người xung quanh.

Bộ trang phục mà cậu mặc chính là bộ khi mà cậu ra mắt lần đầu, với áo lót trong màu đen, áo sơ mi xám mặc trễ vai, bên ngoài mặc áo khoác tím, đeo vòng cổ, diện với quần đen, thắt lưng trắng, một chiếc vòng tay trắng ở cánh tay phải và giày thể thao cao cổ màu đen.

\chrint

``Chuyện gì xảy ra vậy\dots?'' Minato hỏi, đưa tay dụi mắt.

``Biết chết liền đó cậu,'' Kuon nhăn mặt, đáp như thể đang trách nhẹ. ``Bọn tôi cũng vừa mới tỉnh dậy thôi.''

``Trông cậu như thể vừa mới tỉnh dậy sau một trận quẩy điên loạn vậy,'' Kaigou lườm.

``Ặc\dots\ thì đúng là thế mà\dots\ chắc tại tôi uống nhiều quá\dots,'' cậu MC thú nhận, gãi đầu cười trừ.

Suzuki thở dài, giọng vẫn còn yếu: ``\dots Chúng ta\dots\ bị trôi dạt vào đây sao, Onii-sama?''

Cậu Tổng giúp Minato đứng dậy, cũng đang cố định hình lại tinh thần. Nhưng trước khi cậu có thể thở phào, cậu MC đã lại hỏi bằng giọng ngơ ngác: ``\dots Ủa? Ta đang ở đâu đây!?''

``Đến tôi còn không biết nữa đây,'' Kuon đáp, lắc đầu bất lực.

Kaigou khoanh tay, cau mày: ``Tại ai mà ra nông nỗi này, tự biết nhé.''

Chuyện là, vào tối hôm qua, Minato - với vai trò một host sôi nổi và ý tưởng ``nghỉ dưỡng ngoài trời'' đã hăng hái mời mọi người ra biển ``cho có không khí mùa hè''. Dù Kuon - người vừa trở về từ buổi thu âm cho tập phim tài liệu của Fumi - phản đối nhẹ nhàng, Kaigou thì nhìn cậu như nhìn sâu bọ, còn Suzuki thì chẳng dám nói gì, bằng một cách nào đó, cả ba người nhà Hikawa đều bị cậu chàng MC này lôi đi cho bằng được.

Trên du thuyền, Minato bật nhạc, tổ chức tiệc nhẹ và trò chuyện không ngừng. Ngoài với ba thành viên Hikawa (bị thúc ép) ra, cậu còn đi chung với ba thành viên khác. Nhưng rồi, một trận sóng bất ngờ đánh vào mạn tàu, mạnh tới mức cả bốn người bị hất xuống biển. Không một chiếc áo phao, không cứu hộ, và khi họ tỉnh lại, họ đã ở nơi này: một hòn đảo không người, biệt lập giữa đại dương mênh mông.

``Chúng ta là những người duy nhất ở đây sao?'' Minato hỏi, giọng đã tỉnh táo hơn.

``Chắc thế,'' Kaigou đáp, liếc quanh rồi đá nhẹ một vỏ sò gần chân. ``Chẳng thấy ma nào khác ngoài chúng ta cả.''

``Như cô ấy nói đấy,'' Kuon gật đầu. ``Bọn tôi cũng vừa mới tỉnh thôi.''

Cậu MC thở lấy một hơi. ``Vậy thì\dots''

Cậu đứng dậy, chống hông rồi quay sang mọi người: ``Chắc chúng ta cứ đi quanh một vòng nhỉ? Xem có gì ăn, hoặc ít nhất là có cây bóng mát để trú.''

Kuon lặng lẽ gật đầu. Kaigou liếc Minato bằng ánh mắt lạnh như thép, nhưng rồi cũng quay sang phía rừng cây phía xa: ``Tôi nhìn thấy vài bụi rậm đằng kia. Có lẽ vào trong đó xem có gì.''

``\dots Cũng phải. / Đồng ý!'' hai cậu con trai đáp, không một chút phản đối, nhưng với thái độ có phần trái ngược.

Suzuki thận trọng đứng dậy, phủi cát khỏi vai áo. Cô bé không nói gì, chỉ lặng lẽ đi cạnh hai anh chị. Ba người nhà Hikawa, cùng với cậu MC ``tai họa'', bắt đầu tiến vào khu rừng, không ai biết điều gì đang chờ đợi mình phía trước.

\ct

Kaigou bực mình lẩm bẩm khi quan sát người đi trước: ``\vie{\dots Cậu ta đúng là có dáng đi lạ thật.}

Kuon đi cạnh cô, vừa nhăn mặt vừa phủi nhẹ mồ hôi đang lấm tấm trên trán: ``\vie{Trông chẳng khác gì một đứa con nít. Cứ nhảy cà tưng như đang đi dã ngoại ấy.}''

Suzuki thở dài khe khẽ: ``\vie{Lại còn không tỏ ra lo lắng chút nào nữa mới lạ\dots}

Quả thực, trước mắt họ, Fuwa Minato đang hồn nhiên đi giữa rừng cây, vừa nhảy chân sáo vừa ngân nga một đoạn giai điệu nghe rất khó hiểu. Dáng vẻ ấy hoàn toàn đối lập với ba người nhà Hikawa, những người đang bước từng bước cẩn trọng như đang tiến vào lãnh địa của thú dữ.

``\vie{Tôi thậm chí còn thấy ngạc nhiên khi cậu không hoảng loạn trong tình cảnh éo le này đấy\dots} cô chị hất cằm về phía cậu MC, giọng không giấu được vẻ châm biếm.

Sở dĩ họ nói vậy là vì rừng cây mà họ đang đi vào không phải là khu rừng bình thường. Dù là ``giả'', tức là toàn bộ thú hoang bên trong chỉ là hình mô phỏng bằng bìa cứng hoặc giấy ép, nhưng sự bố trí của chúng lại khá rối mắt và có phần rùng rợn, nhất là với ánh sáng hắt xuyên tán lá lốm đốm.

\img{nj-7b}

``Trong này có cả cá sấu\dots\ với gấu ư?'' Minato ngơ ngác nói, nhìn vào một tấm biển dọa dẫm treo lủng lẳng trước mặt.

Suzuki đi sát phía sau, tròn mắt nhìn tấm biển, rồi thì thầm: ``\vie{\dots Chúng là hình tĩnh. Bằng giấy bìa thôi.}''

``\vie{Chắc do kinh phí thấp quá.} Kuon nhún vai, chán nản. \vie{Giống cái hồi hai anh em đi dạo với Saku-chan với Yuika-chan hôm vừa rồi.}''

``\vie{Dù vậy cũng phải cẩn thận,}'' Kaigou nói, tay siết nhẹ cây gậy gỗ cô tiện tay nhặt được từ nãy. Sau đó, cô liếc nhìn một khoảng trống bên cạnh như thể thiếu gì đó, rồi thở dài. ``\vie{Giá mà ông ta ở đây thì tốt.}''

Cậu ngạc nhiên quay sang: ``\vie{Ai cơ?}''

``\vie{Game-san ấy.}''

``\dots À.'' cậu ngẫm một lát, rồi thở ra: ``\vie{Tôi chưa sạc cho ông ta chiều hôm qua, nên đành để ở nhà. Nhưng tôi tưởng bà ghét cay ghét đắng ông ta lắm mà?}''

Cô chị nhún vai, mắt không rời bụi cây phía trước: ``\vie{Ông ta vẫn còn hữu dụng chút ít. Chứ trông chờ vào mấy tên đàn ông xung quanh đây, đúng là\dots}'' cô liếc Minato, rồi thở hắt ra, ``\vie{\dots hên xui lắm.}''

Cô em gái gật đầu, nhẹ giọng đỡ lời cho chị gái: ``\vie{Em thấy\dots\ ít nhất thì Onii-smaa vẫn đỡ hơn đám con trai nơi chị từng sống, đúng không ạ?}''

Kaigou lặng lẽ gật đầu, biểu cảm dịu đi đôi chút. Cô biết rõ, tuy cô vẫn chưa thể tin tưởng hoàn toàn vào phái nam, và vẫn còn mang nhiều vết hằn từ quá khứ, nhưng trong lòng cô biết rõ Kuon không giống như họ. Có điều, sự cảnh giác vẫn luôn nằm đó, bám rễ trong ánh mắt cô mỗi khi nhìn một người đàn ông khác.

Ba người tiếp tục tiến sâu vào bên trong, mắt dán chặt vào từng chuyển động xung quanh. Riêng Minato thì hoàn toàn không để tâm, vẫn đi trước với cái điệu hát nho nhỏ và những cú nhảy vụng về như chú nai vàng ngơ ngác.

Sau hơn mười phút dò đường, khi không thấy gì ngoài bụi cây, lá rụng và vài mẩu côn trùng, cậu Tổng cau mày: ``\vie{Không có gì cả. Chẳng có lối ra, cũng không có cái gì có thể đưa ta rời khỏi đây.}''

``\vie{Xem chừng chúng ta sẽ còn phải mắc kẹt ở đây dài dài\dots}'' cô chị lẩm bẩm, mắt nheo lại.

Bất chợt, cậu MC dừng khựng lại. Đôi mắt cậu mở to, miệng há hốc. Cậu xoay người lại nhìn ba người họ với vẻ mặt như vừa phát hiện ra điều gì kinh thiên động địa: ``Không lẽ\dots\ \textbf{đây là một hòn đảo hoang ư!?}''

Kaigou và Kuon đồng thanh, không cần suy nghĩ: ``Bọn tôi biết từ lâu rồi. / Não cậu chậm quá đó.'' Suzuki thì chỉ im lặng lắc đầu, chống hông nhìn cậu chàng với vẻ nửa chán nản, nửa cảm thông.

Nhưng điều khiến cả ba người bất ngờ không phải là sự nhận ra muộn màng ấy, mà là phản ứng sau đó của Minato. Cậu thở sâu, rồi bất ngờ lấy lại tinh thần rất nhanh, đứng thẳng người dậy, nở nụ cười tự tin và nói to: ``Vậy thì\dots\ chúng ta cần phải sống sót! Và bắt đầu hành trình sinh tồn ngay từ bây giờ!!''

Nói rồi, cậu còn nháy mắt một cái - một cách rất kịch - và cười toe như thể đang đứng giữa trường quay của một chương trình thực tế sinh tồn thực thụ.

Kuon nhìn cậu, ngơ ngác một thoáng rồi buông một chữ, pha lẫn mệt mỏi và sốc nhẹ: ``\dots Bruh,'' rồi bồi thêm trong lúc chống trán: ``Tôi thực sự sốc khi cậu có thể ngay lập tức điềm tĩnh lại như vậy, Fuwa-san.''

Kaigou thì không khách sáo. ``Đừng có mà tỏ vẻ thản nhiên rồi nháy mắt như thế. Kinh quá,'' giọng cô đầy rợn người, toàn thân như nổi da gà.

Suzuki mím môi, cố giấu nụ cười đang trồi lên bên mép, nhưng không nói gì. Với cô, một Minato khùng khùng như thế lại khiến mọi thứ đỡ căng thẳng hơn phần nào, ít nhất là cho tới thời điểm này.

\ct

``Đầu tiên, chúng ta sẽ tìm kiếm thức ăn!'' Minato hô to, đầy khí thế, như thể đang đứng trên sân khấu của chính mình.

Cậu chỉ tay về phía một cái cây mọc nghiêng bên mép cát, trên đó là vài quả thanh long chín đỏ mọng, nằm chễm chệ một cách hoàn toàn phi lý giữa vùng rừng ven biển.

``Cái đó thì công nhận là đúng đấy,'' Kuon đáp, mắt nheo lại, ``nhưng mà tôi thấy có gì đó không đúng\dots''

``Ừ, đúng là có gì đó không ổn thật,'' Kaigou gật đầu, khoanh tay, cau mày nhìn cây trái trước mặt. ``Tôi đâu nhớ loại quả này mọc ở ven biển nhỉ? Tưởng là dừa chứ?''

Suzuki nghiêng đầu, tay chạm nhẹ vào quả thanh long, vẻ trầm ngâm đầy vô tội: ``Em nghĩ đây là giống thanh long quý hiếm nào đó mà chúng ta vừa phát hiện ra-''

*VÚT!*

\img{nj-7c}

Chưa kịp để cô dứt câu, một cái bóng lao tới như gió và cuỗm sạch quả thanh long từ tay Minato. Cả bốn người ngơ ngác, chưa kịp định thần.

Bốn giọng đồng thanh vang lên, người ngạc nhiên, người giật mình, người hoảng sợ:

\begin{dialogue}
    \speak{Minato} N-Này!?
    \speak{Kuon} Cái quái\dots!?
    \speak{Kaigou} Cái gì vậy!?
    \speak{Suzuki} C-Cái gì thế!?
\end{dialogue}

Bóng người đó lăn thêm một đoạn rồi ngồi thụp xuống nền cát, để lộ ra một cô gái với mái tóc xanh thẫm nước biển dài tới lưng gồm hai lọn tóc buộc cao, mỗi bên trên đỉnh đầu đính một chiếc chuông vàng nhỏ kêu leng keng, gợi nhớ đến một chú mèo máy nào đó trong ký ức tuổi thơ. Trên đầu cô là một chỏm tóc ngố vểnh lên như thể chỉ chực bay đi mất.

\chrint

\begin{wrapfigure}[28]{l}{10.5cm}
    \includegraphics[width=10.5cm]{../../../../Image/Book?/kokoro_model.jpg}
\end{wrapfigure}

Amamiya Kokoro (\ruby{天}{あま}\ruby{宮}{みや}こころ, \ruby{A-ma}{Thiên}-\ruby{mi-gia}{Cung}\ Cô-cô-rô). Cô nàng gia nhập vào gia đình Cầu Vồng sớm hơn Minato ba tháng. Cô là một nữ thần (hoặc vu nữ) ở trên trời, có thể giao tiếp với loài rồng. Cô nàng gia nhập vào công ty để có thể hòa hợp với thế giới loài người khi đã trưởng thành. Hiện đang trong quá trình ``học tập và trưởng thành'' ở nhân giới. Cô đặc biệt thích những món ăn có chữ ``long'' trong tên gọi, và tất nhiên, thanh long là một trong số đó.

Cô nổi tiếng với cách phát âm đặc biệt do lưỡi ngắn, thậm chí ngay cả tên ``Amamiya'' của mình cũng được phát âm thành "Amiya-mya" và trở thành biệt danh. Khi đọc tên công ty Nijisanji, cô thường phiên âm thiếu đi một chữ ``ji.'' Đó là điểm cuốn hút của cô nàng.

Với tính cách trẻ con và ngây thơ, Kokoro thường để lại ấn tượng về một hình ảnh rất đáng yêu và trong sáng, mặc dù có thể cô nàng đã rất nhiều tuổi (cô nàng không tiết lộ số tuổi do yêu cầu từ người bố của mình). Khi hứng khởi, cô sẽ bất chợt hát những bài hát ngẫu hứng khó tả hoặc ngân nga.

Do chịu ảnh hưởng từ bộ anime ``Bocchi the Rock!'', cô đã mua một cây đàn guitar điện (Made in Japan Junior Collection Telecaster Rosewood Fingerboard Satin Daphne Blue) và cô nàng đã tập luyện trong hai buổi stream sau đó chỉ để chơi các bài hát tư bộ anime này.

Cô rất thích manga và thường xuyên chia sẻ về chúng trong các buổi phát sóng trực tiếp của mình, đồng thời giới thiệu nhiều tác phẩm manga trên Twitter. Thể loại mà cô ấy thích thường rất đa dạng, chủ yếu là từ thể loại kỳ ảo hay ma thuật.

Bộ trang phục của Kokoro gồm áo sơ-mi trắng bên trong với diềm xếp nếp, bên ngoài trông như một bộ đồng phục học sinh: váy xanh nước biển nhạt được đeo hai dây ở vai có viền đen diềm xếp, bên trái ở hông cài một chiếc nơ cùng màu với chiếc váy. Cô đội một chiếc mũ nồi màu đen và tất dài đến đùi màu trắng, đi giày lười đen. Đằng sau cô còn có một đôi cánh màu trắng.

\chrint

Kokoro ngồi xổm giữa bãi cát, ôm chiến lợi phẩm trên tay và ngấu nghiến nó ngay lập tức, thậm chí không thèm bóc vỏ.

Cả bốn người nhìn nhau, mắt chữ O miệng chữ A.

``Này! \dots Ể!? A-Amamiya!?'' Minato sững sờ.

``Kokoro-nee!?'' Suzuki hét lên, nấp vội đằng sau hai anh chị.

``Kokoro-chan? Tự dưng ở đâu mọc lên vậy?'' Kuon nghiêng đầu, vừa ngạc nhiên vừa bất lực.

``Ơ\dots\ Ê! Từ từ! Em còn chưa-'' Kaigou chưa kịp ngăn lại thì Kokoro đột ngột buông quả thanh long xuống, ôm cổ ho sặc sụa, mắt trợn tròn và mặt đỏ ửng.

*KHỤ KHỤ* *ẶC\dots!*

``\dots bóc vỏ mà! Trời ạ!'' cô chị hét lên, chạy tới, không kịp cứu nguy cho nàng Thiên Cung này.

Kuon lắc đầu thở dài, giọng không giấu được trách móc:  ``Có ai tranh của em đâu mà ăn như thể bị bỏ đói ba ngày vậy\dots''

``K-Kokoro-nee!? C-Chị ổn chứ!'' Suzuki hoảng hốt chạy đến cạnh nàng tiên nữ, khuỵu gối xuống sát mặt cát, nhìn vào khuôn mặt đẫm mồ hôi của cô ấy.

Kokoro chỉ có thể đưa ra ánh mắt tròn xoe ướt nước, rướn người khò khè thều thào một chữ: ``N-Nước\dots!''

\img{nj-7d}

Cậu Tổng lúc này mới bật dậy, sực nhớ ra một điều quan trọng: ``Nhắc mới nhớ, ta sẽ cần nước để sinh tồn thật đấy.''

``Đúng là như vậy!'' Minato gật gù, như vừa được nhắc bài.

``Còn chần chừ gì nữa hả hai người!?'' Kaigou quát. ``\textbf{Mau đi lấy nước cho em ấy ngay!}''

``Nhưng vấn đề là nước ở đâu?''

``Ở trong rừng ban nãy, tôi có thấy một chỗ giống như thác nước. Có thể uống được,'' Kuon nói nhanh.

``Được rồi! Anh sẽ quay lại đó ngay!'' cậu MC đấm nhẹ vào lòng bàn tay, đầy quyết tâm.

Cậu quay lại nhìn Kaigou và Suzuki, gật nhẹ: ``Ở lại canh chừng nhé.''

``Được rồi!'' Suzuki đáp, vẫn còn đang xoa nhẹ lưng Kokoro, trong lúc cô chị thì đang đập mạnh vào lưng của nàng Tiên, người vẫn đang ho sặc sụa.

Ngay sau đó, hai người con trai lập tức quay người, chạy ngược trở lại lối cũ trong rừng rậm để tìm nguồn nước. Những bước chân gấp gáp để lại dấu dài trên cát nóng bỏng.

Và khi bóng dáng họ khuất dần sau tán lá, một giọng nói nhỏ nhẹ như gió thoảng vang lên từ phía sau lưng từ ba người con gái: ``Cứ để đó để cho mình.''

Cô chị giật bắn người, quay đầu lại. Đôi mắt cô trợn lên:
``Khoan đã\dots\ Em là\dots!''

\ct

Hai bóng người khuất dần sau tán cây, vội vã như thể đang tham gia vào một cuộc đua sinh tồn đúng nghĩa.

Kuon đi trước, nét mặt đầy tập trung, chân bước dứt khoát xuyên qua những bụi cây dày đặc. Cậu nhớ rõ vị trí con thác nhỏ vừa thấy ban nãy. Đó là nguồn nước duy nhất mà họ có thể tin tưởng trong hoàn cảnh này.

Minato thì trái ngược hoàn toàn. Cậu vừa đi vừa ngắm cảnh, còn huýt sáo ngân nga một đoạn giai điệu tươi vui như thể đang tham gia chương trình dã ngoại dành cho thiếu nhi. Dù mồ hôi đã lấm tấm trên trán, nét mặt của cậu vẫn giữ nụ cười hồn nhiên kỳ lạ.

Tới nơi, cậu Tổng ngay lập tức quăng ánh mắt về phía vài chiếc chai nhựa cũ nằm vương vãi bên gốc cây gần thác. Cậu nhặt chúng lên, xả nước trước rồi cúi xuống múc nước đầy một cách nhanh chóng, động tác dứt khoát và chính xác.

Còn cậu MC\dots\ sau khi tới nơi, thay vì phụ giúp, lại chậm rãi đi vòng sang bên phải, nơi có một chiếc rương kho báu gỗ cũ kỹ nằm lọt thỏm dưới gốc cây. Cậu mở nắp rương, mỉm cười vui vẻ khi lấy ra một chiếc đầu lâu rỗng, rồi bình thản bước tới thác, múc nước từ chính vật thể có phần đáng ngờ ấy.

``\textbf{Lấy thì lấy nhanh lên đi!}'' Kuon gằn giọng, mắt liếc sang với vẻ khó chịu. ``\textbf{Kẻo em ấy nghẹn mà không chịu nổi đấy!}''

Minato vẫn giữ nụ cười dịu dàng, đáp lại bằng một giọng điệu nhẹ bẫng như mây trời: ``Cứ bình tĩnh. Rồi mọi chuyện lại đâu vào đấy mà\~{}''

Cậu Tổng nghiến răng, thở ra đầy ngao ngán. Lối sống trần tục, vô tư đến phát bực của chàng MC này khiến cậu không thể nào lý giải nổi.

Khi đã gom đầy ba chai nước, Kuon đứng thẳng người, chỉnh lại quai đeo túi rồi quay về hướng bờ biển. Còn Minato thì vẫn vừa đi vừa nhảy nhót theo kiểu rất trẻ con, hai tay xách cái đầu lâu đầy nước, khiến nước sánh ra theo từng bước chân tưng tưng.

``\textbf{Đừng có mà nhảy nữa và đi nhanh lên!!}'' Kuon cuối cùng không chịu nổi, lớn tiếng quát. Cậu suýt thì mất bình tĩnh khi thấy phần nước quý giá đang rò rỉ từng giọt quý như vàng, thứ mà có thể giúp cho cô nàng Tiên đang sắp chết nghẹn ấy.

\ct

``\textbf{Gyaaaaah!!}''

Một tiếng hét chói tai vang lên từ phía Kokoro, người hiện đang bị một ai đó - không phải từ chị em nhà Hikawa - nhấc bổng lên và quay tròn như chong chóng trên không trung.

``N-Này! Em không cần phải làm quá lên như vậy đâu!!'' giọng cô chị đầy bất lực khi cố gắng ngăn cô gái tóc vàng đang cầm chính chân của Kokoro và xoay cô như múa lửa hội làng.

``S-Sara-nee! Em không nghĩ cách đó có hiệu quả đâu!'' cô em gái ở bên cạnh lên tiếng, toàn thân vẫn trốn đằng sau cô chị.

Đúng lúc đó, Kuon xuất hiện, tay cầm ba chai nước đầy, chạy về phía nhóm với hơi thở gấp.

``\textbf{Bọn tôi mang nước về rồi đây!}'' cậu MC tưng tửng đi về phía ba cô gái.

``Khổ ghê cơ, cái cậu này lề mề vãi chưởng-'' cậu Tổng bẹc bội than vãn, rồi điếng người với cảnh trước mặt: ``Ủa? Có chuyện gì vậy?''

``Kuon-kun! Mau tìm cách để em ấy dừng lại đi!'' Kaigou thốt lên, chỉ về phía Kokoro.

Trước mắt họ là một cô gái nổi bật với mái tóc vàng buộc đuôi ngựa lệch phải, trên mái là hai chiếc kẹp tóc hình chữ thập, cùng đôi mắt dị sắc đỏ - vàng sáng lấp lánh. Cô ấy đang dồn hết sức để đẩy quả thanh long khỏi họng của nàng Tiên nữ đen đủi ấy.

\img{nj-7f}

\chrint

\begin{wrapfigure}[30]{r}{11cm}
    \includegraphics[width=11cm]{../../../../Image/Book?/sara_model.png}
\end{wrapfigure}

Người đang quay Kokoro như con quay ấy là Hoshikawa Sara (\ruby{星}{ほし}\ruby{川}{かわ}サラ, \ruby{Hô-shi}{Tinh}-\ruby{ca-oa}{Xuyên}\ Xa-ra), người gia nhập vào Nijisanji cùng quý 4 với Minato. Sara là một cô nàng mang hai dòng máu Anh và Nhật. Cô từng sống ở xứ Anh Quốc ảo lúc cô còn bé, nhưng cô lại dành phần lớn thời gian ở Nhật Bản ảo lâu hơn, nên cô nàng hầu như chỉ nói tiếng Nhật.

Sara được mô tả là một cô gái bướng bỉnh, tinh nghịch, chẳng hạn như không trêu chọc hay đùa giỡn với người xem, vừa cười vừa đọc những đoạn chat như chửi thề và những câu nói tục tĩu. Cô có giọng nói dù hơi trầm, nhưng khá tươi sáng và sống động trong các Liver ở nhà Cầu Vồng.

Sara còn được biết đến với chương trình tư vấn tình cảm trực tuyến của mình - ``Hoshikawa Love Lab'', nơi cô chia sẻ những lời khuyên chân thành về tình yêu và các mối quan hệ cho người xem. Mặc dù còn trẻ, nhưng những lời khuyên của cô luôn được đánh giá cao vì sự chân thành và tinh tế.

Trong cộng đồng VTuber, Sara nổi tiếng với mối quan hệ thân thiết cùng Natsuiro Matsuri từ \hll. Hai người thường xuyên collab và tương tác với nhau trên stream, tạo nên những khoảnh khắc đáng yêu khiến fan hai bên vô cùng thích thú. Họ còn thường xuyên chơi game cùng nhau và có những cuộc trò chuyện thú vị về cuộc sống. Thậm chí một số người còn đùa rằng tính lầy lội và dâm đãng còn ``học của nhau'', dù không ít lần Sara phải đổ mồ hôi hột với cô bạn Lễ hội này.

Ngoài ra, Sara còn được biết đến là một người bạn tốt trong cộng đồng VTuber, luôn sẵn sàng giúp đỡ các đồng nghiệp và tạo không khí vui vẻ trong mỗi buổi collab. Điều này càng làm tăng thêm sự yêu mến của người hâm mộ dành cho cô.

Cô nàng mang hai dòng máu này, trùng hợp thay, có hai màu mắt, với mắt đỏ bên trái và mắt cam bên phải. Sara mặc một chiếc áo ba lỗ màu trắng buộc phía trước, vòng cổ đen, quần soóc ngắn màu xanh, đeo dây đai đùi màu hồng ở bên phải, quần lót cao tới eo (có thể nhìn thấy hai sợi trên eo) màu đen. Cô có đeo dây đai đùi thứ hai ở chân trái dưới có miếng kim loại hình tròn màu vàng, đi tất màu xanh, bốt gấp màu trắng có khéo đằng trước, áo khoác trắng dài tay mặc trễ, có ruy băng màu xanh và một dải ruy băng tóc hai tông màu đỏ cam buộc tóc thành đuôi ngựa lệch sang bên phải.

\chrint

``\textbf{Mau nhả nó ra đi nào!}'' cô nàng hét, tiếp tục xoay Kokoro vòng vòng như đang chơi trò nhào lộn trong rạp xiếc.

*KHẠC!*

Một tiếng phụt vang lên. Từ miệng Kokoro, miếng thanh long mắc nghẹn bị phóng ra không thương tiếc, bay vèo về phía hai chàng trai vừa trở lại.

Kuon thấy thế, phản xạ nhanh như điện giật, liền nghiêng người né qua một bên, tránh được đường bay.

Minato thì không may mắn như vậy. Miếng thanh long đập thẳng vào mặt, khiến cậu ngã ngửa, và bị chính cái đầu lâu rỗng nước đập trúng trán.

\gif{nj-7g}{1}{118}

Cậu Tổng nhíu mày, không giấu nổi mỉa mai: ``\dots Cái giá phải trả cho việc cậu cứ tửng tưng đấy.''

Cậu MC ôm đầu lăn lộn: ``Ôi đau quá\dots''

Kokoro ho sặc sụa, hai tay chống trên bãi cát thở gấp: ``Khụ khụ! M-May quá\dots\ Bà đã cứu tôi rồi\dots''

``Em ổn chứ?'' cô chị lớn vội chạy đến, đỡ nàng tiên nữ đang thở không ra hơi.

Suzuki quay sang, cúi đầu cảm ơn cô gái vừa ra tay cứu giúp: ``Cảm ơn chị nhé, Sara-nee.''

``Khoan\dots\ Sara-chan cũng ở đây sao!?'' cả Kuon và Minato đồng thanh kinh ngạc.

``Hai người quên rồi à?'' Kaigou cau mày. ``Cô ấy cùng với Kokoro-chan và một người khác cũng bị Minato lôi đi với chúng ta mà, nhớ chứ?''

Cô gái tóc vàng quay lại, nở một nụ cười nhàn nhạt: ``Đúng đó. Lúc tôi tỉnh lại thì đã ở đây rồi. Chỉ nhớ rằng lúc tôi đang lướt điện thoại thì bỗng dưng cả con thuyền rơi ụp xuống biển rồi khiến tôi trôi về đây thôi.''

``Vậy nghĩa là\dots\ vẫn còn một người nữa sao, Sara-nee?'' cô em gái hỏi.

Cô nàng xứ Anh nhún vai: ``Có thể nói là thế. Mà\dots\ đây là đâu vậy?''

``Biết chết liền á. Chỉ biết đây là một hòn đảo hoang.''

``\textbf{Đây là\dots\ một hòn đảo hoang!!}''

\begin{wrapfigure}[29]{l}{8cm}
    \includegraphics[width=8cm]{../../../../Image/Book?/kakeru_model.jpg}
\end{wrapfigure}

Một giọng nói nam trầm rõ ràng vang lên từ phía sau nhóm năm người đang tụ tập bên bãi cát. Tất cả họ quay đầu lại và trông thấy một cậu chàng cao ráo với mái tóc đen, mắt đỏ, diện bộ trang phục thường thấy. Trên cổ còn đeo một chiếc tai nghe màu đen-xám, như một nhạc sĩ luôn sẵn sàng thu âm mọi lúc.

Minato, Kokoro và Sara đồng thanh kêu lên trong ngỡ ngàng: ``\textbf{Yumeoi!?}''

\chrint

Yumeoi Kakeru (\ruby{夢}{ゆめ}\ruby{追}{おい}\ruby{翔}{かける}, \ruby{Giu-mê}{Mộng}-\ruby{oi}{Truy}\ \ruby{Ca-kê-rư}{Tường}), nhà soạn nhạc - ca sĩ kỳ cựu của Nijisanji, một trong những người anh cả của dàn nam nhà Cầu Vồng, và lớn hơn Kuon bốn tuổi. Cậu gia nhập nhánh SEEDs từ tận năm 2018, tức là trước cả ba người họ khoảng một năm.

Một nhà soạn nhạc - ca sĩ kỳ cựu của nhà Cầu Vồng, một người thích sáng tác âm nhạc và hát vang những ca khúc. Ước mơ của cậu là muốn được mọi người trên toàn thế giới thưởng thức những kiệt tác âm nhạc của mình.

Tuy nhiên, đôi khi anh cũng có những lúc hoàn toàn trái ngược. Với các fan của mình, tiếng cười thật sự của Kakeru khá đáng sợ và mang lại cảm giác không mấy tốt đẹp. Điều này có lẽ bị ảnh hưởng từ buổi collab Hayato, nơi anh bị cho là có xu hướng sát nhân.

Kakeru thích trêu chọc đồng nghiệp và là một người rất hòa đồng, khiến anh trở thành một ứng cử viên sáng giá cho vị trí MC trong các sự kiện. Anh thực sự xuất sắc trong vai trò điều phối viên, và được khen ngợi rất nhiều về khả năng làm việc ổn định với Yumeo trong bất kỳ sự kết hợp nào với các thành viên khác.

Mặc dù được biết đến là một ca sĩ-nhạc sĩ, nhưng anh lại sống trong một live house vì không có nhà và thường xuyên ăn giá đỗ cho bữa ăn vì không kiếm được nhiều tiền từ công việc chính.

Một điều thú vị là mối quan hệ giữa anh và Hoshikawa Sara - một đàn em trong công ty. Ngay từ cuộc gọi điện đầu tiên và buổi collab đầu tiên, Sara đã gọi anh là "papa" và khiến anh trở thành một người cha trong chớp mắt. Anh được biết đến với biệt danh "Sleeveless Psycho Papa", và sau đó họ thực sự được coi như cha con thật sự. Mối quan hệ này được gọi là Gia đình Yume Hoshi (夢星家), dù tất nhiên, anh vẫn luôn phủ nhận điều này.

Khả năng tiếng Anh của Kakeru khá tốt. Tuy nhiên, do cảm thấy áp lực, anh thường chỉ nói những câu đơn giản. So với Kuon, Kaigou hay thậm chí là Suzuki - những người thường xuyên sử dụng cả ba ngôn ngữ Etsunan, Nhật và Anh một cách tự nhiên - có lẽ đây là một điểm yếu của anh.

Bộ trang phục thường thấy của anh bao gồm áo khoác trùm đầu màu trắng không tay, một chiếc áo đỏ không tay mặc bên trong, quần đen với phần ống có thể nhìn thấy, bốt đen dây đỏ, vòng tay đen ở tay phải và đeo tai nghe quanh cổ.

\chrint

Kaigou thì lại khoanh tay, chép miệng: ``Duh, tất nhiên chúng ta đang ở trên đảo hoang rồi. Anh nghĩ đây là đâu vậy hả, Kakeru?''

Kuon cũng nghiêng đầu hỏi, vừa bối rối vừa nghi hoặc: ``Với cả, anh đang làm gì ở đây vậy, Kakeru-san?''

Cậu chàng đeo tai nghe kia vẫn đứng chống hông, giọng oang oang như thể đang phát biểu tại sân khấu hòa nhạc: ``Đã lâu lắm rồi tôi mới có thể gặp được nhiều người như vậy đấy! Phải mười năm rồi chứ ít gì!''

Ngay lập tức, mồ hôi lạnh rịn ra trên trán từng người một.

Minato nhăn mặt: ``Chắc chắn là không lâu đến thế ạ\dots''

Cậu Tổng thì thản nhiên: ``Em còn chưa gặp anh lần nào.''

Khuôn mặt Yumeoi lại bừng sáng ngay sau đó như thể lời trêu đùa nhẹ nhàng của hai người kia chẳng ảnh hưởng gì đến tâm trạng anh ta. Anh giơ nắm tay lên, cười rạng rỡ: ``Dù gì thì, anh rất mừng khi gặp mọi người ở đây! Giờ anh có thể về nhà được rồi!''

Một thoáng im lặng bao trùm bãi biển. Mặt ai nấy đều sụp xuống, như thể tia hi vọng vừa lóe lên đã tắt phụt ngay sau đó.

Suzuki nhỏ tiếng nói: ``Vấn đề là\dots\ ta về như thế nào đây ạ?''

Minato nhún vai: ``Bọn em cũng bị lạc y như anh thôi.''

Hai câu nói cứng như đá tảng đập thẳng vào lòng Kakeru, khiến anh đơ ra như tượng phật đá cổ, hai mắt mở lớn. Sau một thoáng lặng người, anh quỵ xuống cát, thì thầm mấy từ nghe không rõ: ``Không\dots\ thể\dots\ nào\dots''

Kaigou nhún vai, thở dài. ``Dù có muốn hay không thì chúng ta cũng phải chấp nhận thôi,'' rồi quay sang nhìn Kuon, nhấn mạnh: ``Ngoài trừ Kuon-kun ra, còn lại mình toàn gặp phải mấy tên dở người\dots''

Đứng bên cạnh chị gái, Suzuki tiến lên vài bước, nghiêng đầu hỏi với vẻ lịch sự: ``A\dots\ Kakeru-san\dots\ chú bị lạc vào đây bằng cách nào vậy ạ?''

Yumeoi bừng tỉnh lại, ngước lên từ chỗ ngồi và đáp bằng vẻ thành thật đến ngô nghê: ``Chú gì? Anh thôi! Chả là, anh đang\dots\ núp dưới tầng hầm của con tàu để tập trung viết nhạc.''

``Tàu\dots? Ý anh là con tàu của Fuwa-nii sao?'' cô bé quay sang phía Minato như thể muốn xác nhận.

Cậu MC gật đầu, hơi ngơ ngác: ``Hình như đúng là vậy thật\dots\ Chẳng trách tại sao anh thấy có cái gì lù lù ở dưới\dots''

Suzuki quay lại, nhíu mày nhẹ: ``Con tàu đó bị chìm giữa chừng do sóng lớn mà\dots\ vậy nghĩa là anh cũng bị cuốn trôi tới đây?''

``Có vẻ là vậy\dots'' cậu Nhạc sĩ cười gượng, gãi đầu, ``Lúc anh tỉnh dậy thì đã ở đây rồi.''

Sara khoanh tay, nhướng mày nhìn Minato: ``Anh đưa người ra biển tiệc tùng mà không kiểm tra kỹ tàu hả?''

Minato vội vàng chống chế: ``Tôi tưởng mọi người xuống hết rồi chứ\dots''

Kaigou lườm Minato bằng ánh mắt nửa khinh nửa mệt: ``Tổ chức tiệc gì mà như buôn lậu người không bằng\dots''

Lúc này, Kuon đập nhẹ hai tay vào nhau, kéo tất cả trở lại vấn đề chính: ``Thôi được rồi. Có gì tranh cãi thì cứ để sau đi. Vậy bước tiếp theo trong công cuộc sinh tồn trên hòn đảo này là gì đây?''

Tất cả mọi người đồng loạt nhìn nhau, ánh mắt lộ rõ sự\dots\ vô định. Dường như không có ai ở đây biết được họ sẽ làm gì tiếp theo cả.

\uline{Vài tiếng sau.}

Kể từ sau cuộc hội ngộ đầy bất ngờ của sáu người ấy, cả nhóm đã buộc phải tìm cách sinh tồn trên hòn đảo hoang vu, nơi không có bất kỳ dấu hiệu nào của con người, và càng không có tín hiệu để gọi cứu trợ.

Minato với tư cách là một dẫn chương trình, đã đứng ra chia nhóm, cố gắng ``làm chủ sân khấu'' như thường lệ. Cậu giơ tay, hô to: ``Vậy thì ta sẽ chia nhóm như sau! Tôi, Kokoro và Sara sẽ đi kiếm thức ăn! Kakeru sẽ\dots\ ờm, đứng gác! Kuon tìm chỗ nước uống! Kaigou và Suzuki sẽ\dots\ ơm, dựng sân khấu!''

Song, kế hoạch ấy đã ngay lập tức gây ra nhiều ý kiến trái chiều, nhất là từ phía chính các thành viên thuộc nhà Cầu Vồng.

Cô nàng xứ Anh nhíu mày: ``Tại sao lại là em đi với anh?''

Suzuki nhẹ giọng nhưng cũng nêu ý kiến của mình rõ ràng hơn: ``Nếu Onii-sama và Onee-sama lo lửa, thì em nên ở cùng họ để hỗ trợ ạ.''

``Cái sân khấu thì giúp ích được gì cho việc sinh tồn hả Minato?'' cô chị lớn thì cau mày, lườm nguýt cậu MC khiến mọi người vừa gật gù đồng tình vừa đổ mồ hôi hột.

Cậu Ca sĩ thì ngẩn tò te: ``Tôi\dots\ gác cái gì cơ?''

Nhìn cảnh hỗn loạn trước mặt, Kuon chỉ biết thở dài, sau đó tiến lên và sắp xếp lại toàn bộ đội hình theo cách mà ai cũng thấy hợp lý hơn, với tư duy của một Tổng quản. ``Minato-san sẽ đi với Kokoro-chan để hái trái cây. Sara-chan và Kakeru-san tìm nguyên liệu hoặc đồ hữu dụng dọc bờ biển. Ba người nhà Hikawa bọn tôi sẽ lo phần nhóm lửa và chuẩn bị chỗ nghỉ. Mọi người thấy thế nào?''

Minato nhìn Kuon với ánh mắt như thể bị ``cướp sóng truyền hình'', nhưng nhìn quanh thấy không ai phản đối, đành gãi đầu ngậm ngùi: ``\dots Thì\dots\ thôi cũng được vậy.''

Dù sao thì, phần lớn thời gian bốn thành viên nhà Cầu Vồng vẫn mải chơi đùa hơn là nghĩ đến chuyện sinh tồn hay tìm đường về đất liền. Trái ngược lại, ba người nhà Hikawa - Kuon, Kaigou và Suzuki - lại tập trung hết sức cho các công việc cần thiết để tồn tại qua đêm.

\ct

Khi buổi trưa đến, ba người Hikawa tụ họp ở một khoảng đất trống bên bìa rừng để bắt đầu nhóm lửa.

Kaigou ngồi xổm xuống, tay cầm hai que củi: ``\vie{Này Kuon-kun, cậu lấy hộ tôi mấy thanh củi trên kia được không?}''

``\vie{Rồi rồi\dots\ Cô định nhóm lửa thế nào đây?}''

``\vie{Chắc là sẽ dùng lực ma sát thôi nhỉ?}''

``\vie{Nói thì dễ lắm\dots}''

Bên cạnh hai người, Suzuki nhanh nhẹn lấy một miếng gỗ lớn lót dưới đáy, giữ cho ổ nhóm lửa khô ráo. Cô cẩn thận đặt các mảnh vỏ dừa khô mà mình tìm được vào giữa, nói nhỏ: ``Em nghĩ nếu thêm ít lá khô thì sẽ bén nhanh hơn đấy ạ.''

``\vie{Ừm, được đấy Suzu!}'' Kaigou gật đầu, cười nhẹ.

Trong khi ba người đang loay hoay thổi từng đốm lửa đầu tiên, bốn người còn lại đang chơi trò đập dưa hấu một cách đầy hỗn loạn, hoàn toàn bỏ bê công việc.

\img{nj-7j1}

Sara đang bịt mắt, tay cầm cây gậy, hét lên hùng hổ: ``\textbf{Dưa hấu ở đây đúng không nào!? Yaaaah!!}'' và lao tới.

Kokoro thì quýnh quáng vừa chạy vừa la thất thanh: ``Không! Đừng mà! Mình không phải là dưa hấu đâu!!''

Minato thì luống cuống chỉ vào quả dưa hấu: ``Sara-chan! Quả dưa hấu đang ở dưới chân em kìa! Ở dưới!''

Chỉ có Kakeru là tự nhiên như ở nhà, khi cậu lặng lẽ nhặt vài chiếc vỏ sò mà không để tâm tới cuộc rượt đuổi phía sau.

Nhìn cảnh tượng đó từ xa, Kuon chỉ biết thở dài ngán ngẩm: ``Bọn mình đang cố tìm cách sinh tồn, còn bọn họ thì đang đi nghỉ mát thế kia\dots''

\ct

Sự việc tưởng chừng còn tiếp tục điên rồ hơn khi Minato bị cá mập dí giữa biển trong lúc đang đi bơi.

\img{nj-7j2}

``\textbf{Ối dồi ôi! Bớ người ta có cá mập!!}'' mặt cấu xám hẳn đi, loạng choạng bơi vào bở.

Kokoro ở trên cạn hét như thúc giục: ``\textbf{Mau bơi vào đi, Fuwa-san!!}''

Kaigou thì chạy ra nhìn cảnh tượng trước mắt, đổ mồ hôi hột: ``Tôi không nghĩ là vùng biển này có nhiều cá mập đến thế\dots''

Thấy rằng cậu bạn bị bơ đẹp, cậu hét thật lớn, nhấn mạnh từng chữ như một diễn viên quần chúng đang tuyệt vọng lên vai chính: ``Dù, sao, thì! Mau cứu tôi đi chứ không tôi thành sashimi mất!!''

Cậu Tổng cầm sẵn chiếc phao, hét vọng: ``\textbf{Bám lấy cái phao đi!}''

Suzuki chạy theo anh trai, cầm lấy sợi dây buộc quanh phao, hỗ trợ kéo từ xa: ``Onii-sama, em giữ dây cho! Anh mau kéo Fuwa-san lên!'' Cùng lúc đó, đích thân Kaigou xuống dưới biển để đuổi hai con cá mập đi.

Cuối cùng, sau màn hú hồn hú vía, Minato được đưa vào bờ an toàn, tuy rằng ướt như chuột lột.

Kakeru vẫn tỉnh bơ, giọng đều đều: ``Mà tôi tưởng cá mập sợ tiếng động chứ?''

Sara nhăn mặt: ``Chắc nó thấy anh Minato-san bơi kiểu kỳ quặc nên tò mò thôi.''

\ct

Buổi chiều đã buông xuống, khi ánh mặt trời dần lặn, cả nhóm ngồi xung quanh đống lửa mà ba người nhà Hikawa đã nhóm từ chiều. Bên cạnh họ, một dòng chữ to ``S.O.S'' được viết bằng đá và gỗ lên bãi cát.

\img{nj-7j3}

Cô nàng xứ Anh đập hai tay phủi bụi: ``Và thế là\dots\ xong!''

Kakeru vươn vai, chắp hai tay sau gáy: ``Giờ ta chỉ còn cách chờ thôi nhỉ?''

Kokoro thì rơm rớm lệ: ``Nhưng\dots\ ta sẽ phải chờ đến khi nào đây?''

Kaigou nhìn về phía chân trời, lặng im một lúc rồi nói: ``\dots Chị cũng không biết nữa.''

Minato nghe câu nói chán nản của cô chị, gục đầu xuống gối: ``Xem chừng ca này khoai rồi đây\dots''

Dù không ai nói ra, nhưng trong lòng họ đều hiểu: nơi đây sẽ là ``ngôi nhà tạm thời'' của họ, ít nhất là cho đến khi điều kỳ diệu xảy đến.

\uline{Đêm tối hôm đó.}

720 \ruby{giờ}{phút} kể từ khi nhóm bảy người ấy bị trôi dạt vào hòn đảo hoang giữa biển trời. Với họ, từng phút từng giây trải qua dưới cái nắng gắt, gió biển và những tiếng la hét thì cũng dài không kém gì một tháng.

``Mình thật sự đã ngán hoa quả với cá lắm rồi\dots'' Sara thở dài đầy tuyệt vọng, buông rơi nửa miếng thanh long trên tay.

``Thì còn gì ngoài thanh long với cá đâu,'' Kuon gắt nhẹ, cũng chẳng khá hơn. ``Còn hơn là ăn cát.''

``Nhưng mà\dots\ chỗ lương thực này chắc mai là cạn sạch,'' Kaigou thở ra, nằm xoài dưới bóng lá dừa, ``Không biết còn chịu được bao lâu nữa\dots''

``Chắc là\dots\ ta sẽ chẳng bao giờ thoát khỏi nơi này nhỉ\dots'' Kuon lẩm bẩm, mắt nhìn trân trân lên trời đêm.

Tưởng như mọi tia hy vọng đều đã vụt tắt, thì bất ngờ từ phía bên kia rừng cây vang lên tiếng gọi hớn hở của Kokoro: ``\textbf{Mọi người ơi! Lại đây xem này!!!}''

Suzuki cũng chạy theo từ trong bóng tối, thở hổn hển: ``\textbf{Thật đó! Em và Kokoro-nee phát hiện ra một thứ lạ lắm!}''

``Sao thế? Gì mà ồn ào thế\dots'' cô chị lật mình dậy, cằn nhằn.

``Cứ lại đây đi! Nhanh lên!'' giọng của nàng Tiên vang vọng từ trong rừng.

``Thở dài\dots\ Biết vậy,'' cậu Tổng đứng dậy, phủi cát trên người, ``Kaigou-chan, gọi Minato với Kakeru-san dậy đi. Chắc không phải là một con cá biết bay đâu nhỉ?''

Cô chị lớn gật đầu miễn cưỡng rồi dùng chân đá nhẹ vào cậu MC và cậu Nhạc sĩ đang ngủ ngáy gần đó: ``Này, Minato. Yumeoi. Lại có `phát hiện' nữa này.''

``Ứ\dots'' Minato rên rỉ, ``Giờ này mà cũng bắt đi khảo cổ à\dots''

``Cứ đi đi rồi biết.'' Kaigou khoác vai kéo cả hai dậy, không cho họ cơ hội để thức dậy đàng hoàng.

\img{nj-7k}

Cả nhóm bảy người bước vào rừng, băng qua cái thác nước quen thuộc, nơi mà Kuon từng phải hét vào tai Minato vì cậu ta nhảy tưng tưng khi đi lấy nước. Lũ gấu, rắn và cá sấu giấy vẫn nằm đó như những vật thể trang trí lạc đề.

``Lần này là gì nữa đây?'' cậu Tổng hỏi, mắt mỏi nhừ vì thiếu ngủ.

``Là\dots\ cái mà em tình cờ thấy được khi đuổi theo con bướm!'' nàng Tiên nữ hớn hở.

``Em còn rảnh hơn cả mình tưởng,'' cậu Ca sĩ lẩm bẩm.

``Em tìm thấy thuyền à?'' cô chị lớn hỏi.

``Còn hơn thế nữa cơ!'' cô em thêm vào, ``Bọn em đã men theo đường mòn phía bên kia và rồ\dots''

Họ không còn cách nào ngoài việc đi theo hai cô gái nhỏ ấy. Khi họ tới được phía bên kia của hòn đảo\dots

``Cái quái gì\dots?'' Kuon đứng khựng lại.

Trước mắt họ là thành phố Cầu Vồng, hiện ra rõ ràng, lấp lánh ánh đèn ở ngay phía bên kia bãi đá, với tàu thuyền đi lại nhộn nhịp. Không phải vài hải lý xa xôi. Không phải ``trôi dạt tận chân trời góc bể.''

Nó chỉ cách họ đúng một khoảng tương ứng với một con hồ.

\img{nj-7l}

``Đùa nhau à.'' Kuon và Sara đồng thanh.

``Nó\dots\ gần đến thế ư!? \textbf{A-Anh đã làm gì trong suốt mười năm vậy?}'' Kakeru hét lên, như thể bị phản bội bởi chính thế giới này.

``Đừng có phóng đại nữa!'' Minato gắt, ``Mới có nửa ngày thôi!''

``Thế mà bọn mình đi lòng vòng trên đảo cả buổi, dựng lều, nhóm lửa, đập dưa\dots'' Kaigou vò đầu, thở dài chán chường, ``Đúng là trò hề tập thể.''

Không còn thời gian để lăn tăn. Cả bọn buộc phải gom chút sức lực cuối cùng bơi về đất liền.

Cảnh tượng trở nên hỗn loạn, không kém gì một nhóm người đang cố gắng trốn chạy khỏi zombie.

\img{nj-7m}

Sara đang cưỡi lưng Kakeru, vừa đánh bôm bốp lên đầu anh vừa la: ``''\textbf{Papa là đồ ngốc! Đại ngốc! Đại đại ngốc!}''

``Đừng đánh đầu tôi nữa! Tôi còn phải sáng tác đấy!!''

Kaigou thì cõng Kuon một cách bất đắc dĩ (cậu Tổng không biết bơi): ``Bám chặt vào! Không là tôi bỏ cậu lại luôn đấy!''

``Chỉ cần không chết đuối là được\dots'' cậu rên rỉ.

Minato trôi theo một cái thùng phao, miệng cười toe toét, còn Suzuki thì kéo tay Kokoro, vừa đẩy vừa động viên: ``Kokoro-nee cố lên! Em không bỏ chị lại đâu!''

Đằng xa, vỏ sò mà Kakeru nhặt hồi sáng lặng lẽ trôi theo con sóng, khắc dòng chữ nguệch ngoạc bằng đá: ``720 phút. 7 người. 0 lý do hợp lý để bị mắc kẹt!''

%==========================================TRUYỆN PHỤ========================================
\omk

Khi tất cả đã leo được lên bờ, ai nấy đều ướt như chuột lột, mệt lả vì kiệt sức, tay chân nhũn ra như bún sau hơn mười hai tiếng ``sinh tồn''. Họ đều nằm phịch xuống bãi cát của vùng ngoại ô thành phố Cầu Vồng. Phía xa xa, ánh đèn đường đang nhấp nháy như mời gọi. Âm thanh ô tô, tàu điện và tiếng loa phát thanh quen thuộc vang lên, báo hiệu rằng họ đã trở lại với nền văn minh.

Sara vẫn còn đang đập bồm bộp lên lưng ``ông bố bất đắc dĩ'' Kakeru: ``Ai lại để con gái cõng đi thế hả? Từ giờ con sẽ không mua đĩa nhạc của papa nữa!''

``Thì ai bảo em xông vào nước trước làm gì!'' Kakeru vừa ôm lưng vừa càu nhàu, mặt đỏ bừng vì xấu hổ.

Cách đó không xa, Kaigou chống tay thở dốc, vừa cõng Kuon vừa lẩm bẩm: ``Biết thế hôm trước tôi luyện tập thể lực nhiều hơn chút là đỡ khổ rồi\dots''

Kuon, ngồi phệt xuống bên cạnh, vẫn đang rũ rượi: ``Lần sau mà ai nói đi dã ngoại kiểu mới thì tôi xin kiếu.''

Suzuki, dù nhỏ tuổi nhất nhưng có vẻ hồi phục nhanh nhất. Cô bé đã đứng dậy, phủi cát trên váy và giúp Kokoro vắt bớt nước ra khỏi mái tóc dài xanh rì của cô nàng. ``Kokoro-nee, chị đỡ hơn chưa?''

``Uhm\dots\ vẫn hơi lạnh một chút\dots'' Kokoro vừa nói vừa hắt xì, mũi đỏ bừng. ``Đáng ra chị không nên ăn nhiều thanh long như vậy\dots''

Minato, nằm dang tay trên cát như ngôi sao, lẩm bẩm một cách\dots\ triết lý: ``Thật lạ\dots\ chỉ cách nhau một eo biển nhỏ mà cảm giác như sống hai cuộc đời vậy\dots''

Không ai đáp lại, nhưng nét cười mệt mỏi của tất cả họ như một sự đồng tình thầm lặng. Họ vừa thoát khỏi một cuộc ``phiêu lưu đảo hoang'' mà ban đầu tưởng như không có hồi kết, và giờ đây, điều quan trọng nhất là họ vẫn ở bên nhau.

Một cơn gió đêm lướt qua, mang theo mùi muối biển pha trộn với mùi khói lửa còn vương lại trên áo hai chị em nhà Hikawa.

``Thôi thì, mai hẵng tính tiếp,'' Kuon nói, khẽ ngửa đầu nhìn bầu trời đầy sao. ``Cũng đâu phải ngày nào ta cũng có được một chuyến sinh tồn đáng nhớ như thế\dots''

Đối với Minato, cậu vẫn không thể nào quên được trải nghiệm bị hai con cá mập dí dưới biển. Dòng tự sự của cậu trong cuốn nhật ký có ghi:

\txtbx{\hfill ``Fuwa Minato, ngày 16 tháng Tám (thứ Hai)\\
    Hôm nay, tôi thấy mình đang bị mắc kẹt trên một hòn đảo hoang. Tôi đã bị săn đuổi bởi cá mập và suýt chút nữa đã bỏ mặc tại nơi này.\\
    Tôi rất mừng khi tôi đã có thể trở về đất liền an toàn.''}

Và dường như đó cũng là cảm xúc chung của tất cả mọi người trong những ngày đó, cũng có thể sẽ trở thành những ký ức đáng quên nhất trong đời họ, vì chưa bao giờ sự việc đáng ra lại có thể giải quyết một cách ngắn gọn hơn như thế.

Ở đâu đó ngoài khơi, con thuyền gỗ mỏng manh mà Minato từng chèo giờ đang trôi dạt vô định. Bên trong khoang dưới, một cái vỏ sò nhỏ khác mà Kakeru cất kỹ - chiến lợi phẩm duy nhất anh mang về - lăn lóc theo từng con sóng nhẹ. Trên thân vỏ sò ấy, có khắc dòng chữ nguệch ngoạc bằng dao nhỏ: ``720 \sout{giờ} phút sống sót cùng lũ điên.''
\end{document}